% Level 2, Week 4 -- While Loops and Game Loops
\lesson{4}{While Loops and Game Loops}{A while loop repeats until something changes, and every game is a loop.}
\lessoninfo{Concepts: while (I), break (I), input validation (I), nested loops for pixel art (E)}{Tools: Thonny, paper, pencils}{CSTA: 2-AP-12, 2-AP-16, 2-IC-23}

\begin{objectives}
  \item The student can write a \texttt{while} loop that stops because a variable changes, and explain why a loop that never changes runs for ever. (L2.G2; L2.W04.1)
  \item The student can use \texttt{while True} with \texttt{break} or \texttt{return} to keep asking until an answer is valid. (L2.G2; L2.W04.2)
  \item The student can use a loop inside a loop to visit every pixel of a grid made from a list of strings. (L2.G7; L2.W04.3)
  \item The student can say what makes a game fair and name one trick a game should not play on its players. (L2.G9; L2.W04.4)
\end{objectives}

\begin{materials}
  \item Print (one per student): workbook chapter 4, homework 4.
  \item For the warm-up: a pencil and a scrap of paper per student (for the secret number). Nothing else.
  \item On each Mac: in Thonny press \keys{Cmd}+\keys{O} and open \texttt{starter.py} in the folder \texttt{2-intermediate/} \texttt{exercises/week04/ex-guess-number}.
  \item Know where Thonny's red \menu{Stop} button is (top toolbar, a red stop sign). Today someone will need it!
  \item A timer.
\end{materials}

\needspace{16\baselineskip}
\section{At a glance}
Times are a guide: set a timer for each phase and move on when it rings.
\begin{timeplan}
  \timeslot{0--8}{Warm-up}{Unplugged \enquote{Guess my number} with higher/lower. Count the guesses.}
  \timeslot{8--18}{Concept}{\texttt{for} versus \texttt{while}. The loop condition. What changes inside the loop? The infinite loop and the Stop button.}
  \timeslot{18--28}{Guided demo}{Run the Guess My Number starter: it never ends. Fix it live with \texttt{tries} and \texttt{break}. Then type \enquote{banana}.}
  \timeslot{28--52}{Hands-on build}{Pairs: Exercise 4.1 (input validation), then 4.2 Pixel Art. Swap every 8 minutes. Last 6 minutes: design your own pixel picture.}
  \timeslot{52--60}{Wrap-up}{Fair-play talk, exit ticket, reflection, hand out homework 4.}
\end{timeplan}

\section{Warm-up (0--8 min)}
\begin{unplugged}[Guess my number]
You think of a secret number from 1 to 20 and write it on a scrap of paper (face down). Students take turns guessing.
After each guess you say only \enquote{higher}, \enquote{lower} or \enquote{correct!}. One student counts the guesses with tally marks.
Play twice, then let a student be the secret-keeper.
\par\textbf{Ask:} \enquote{What did I do after every guess?} (Checked it and answered.) \enquote{When did the game stop?} (When the guess was right.)
\enquote{What is the smartest first guess?} (10: the middle cuts the numbers in half. With halving, 5 guesses are always enough for 1 to 20.)
\par\textbf{The point:} the game is a loop. \emph{While} the guess is not correct: guess, check, answer. That is today's new loop.
\end{unplugged}

\section{Concept (8--18 min)}
\begin{teacherscript}[Say]
\enquote{In Level 1 you used \texttt{for} loops. A \texttt{for} loop knows in advance how many times to go round: \texttt{for i in range(4)} is
always four times. But in our guessing game nobody knew how many guesses it would take. For that we need a
\textbf{while loop}. It keeps going round \emph{while} something is true, and checks again at the top every time.
Important: something inside the loop must change, or the check never becomes false and the loop never stops.
That is called an \textbf{infinite loop}. If it happens, press the red Stop button. It is not a disaster: every programmer does it.}
\end{teacherscript}
Write this on paper or a whiteboard and trace it together, writing the value of \texttt{n} each time round:
\begin{pycode}[Countdown with while]
n = 3
while n > 0:
    print(n)
    n = n - 1        # the change that lets the loop end
print("Go!")
\end{pycode}
Prints \texttt{3}, \texttt{2}, \texttt{1}, \texttt{Go!}. Ask: \enquote{What happens if I rub out the line \texttt{n = n - 1}?} (It prints 3 for ever.)
Then show the other pattern: \texttt{while True:} goes round for ever on purpose, and \texttt{break} jumps out of the loop when we say so.

\begin{analogy}[The car journey]
\enquote{Are we there yet?} is a while loop. While we are not there: drive a bit, then ask again. It works because the car moves.
If the car stays parked, you can ask \enquote{Are we there yet?} for ever and the answer never changes: an infinite loop.
\texttt{break} is like the driver saying \enquote{We are here, everybody out!}
\end{analogy}

\section{Guided demo (18--28 min)}
\begin{workedexample}[Fixing the game loop together]
\begin{enumerate}[leftmargin=1.4em]
  \item Open \texttt{week04/ex-guess-number/starter.py}. Press \menu{Run} (\keys{F5}). Guess until you hit the secret number.
        It says \texttt{Correct}\ldots{} and asks again! Press the red \menu{Stop} button. Say: \enquote{An infinite loop. Why?}
  \item Point at \texttt{while True:} in \texttt{play()}. \enquote{Nothing in here ever leaves the loop.} Under \texttt{print(result)}, at the same indent, type
        \texttt{if result == "Correct":} and, indented one more step on the next line, \texttt{break}.
  \item Run again. It stops, but says \enquote{You got it in 0 tries}. Ask the students for the fix: \texttt{tries = tries + 1} inside the loop (TODO 1).
  \item Run again and type \texttt{banana} as a guess. \texttt{ValueError} (read bottom-up: \texttt{invalid literal for int()}). Leave that one for the students: it is TODO 2.
\end{enumerate}
\end{workedexample}
The finished program, for you. Note how \texttt{ask\_number()} uses \texttt{while True} and leaves with \texttt{return} only when the answer is good:
\pyfile{../exercises/week04/ex-guess-number/solution.py}

\section{Hands-on build (28--52 min)}
\begin{pairroles}[Driver, navigator, tester]
Swap roles when the 8-minute timer rings (about minutes 36 and 44).
\textbf{1 student:} you are the navigator and the tester: type silly answers like \texttt{banana}, \texttt{99} and \texttt{-3} to try to break the game.
\textbf{2 students:} driver and navigator. \textbf{3 students:} driver, navigator and \textbf{tester} (the tester plays the game after every
change and tries to break it). \textbf{4 students:} two pairs; at minute 44 each pair tries to break the other pair's guessing game.
\end{pairroles}

\begin{enumerate}[leftmargin=1.4em]
  \item \textbf{Exercise 4.1 Guess My Number} (about 10 minutes): finish TODO 2, input validation. Hint for stuck pairs:
        \texttt{answer.isdigit()} is \texttt{True} only when every character is a digit, so \texttt{int(answer)} is then safe.
  \item \textbf{Exercise 4.2 Pixel Art} (folder \texttt{ex-pixel-art}, about 10 minutes): only the top row is drawn.
        Fix: wrap the drawing in the outer loop \texttt{for row in grid:}, and after each row move down
        with \texttt{y = y - size}. Last, students finish the function \texttt{count\_pixels()} (TODO 2).
  \item After each exercise: open \texttt{check\_my\_work.py} and press \menu{Run}.
\end{enumerate}

\begin{bughunt}[The loop that never ends]
In the starter, the game loop has no way out. Ask the pair to say in words: \enquote{The loop stops when \ldots{}} If they cannot finish the
sentence, the loop has no exit. Two questions find almost every infinite loop:
(1) \enquote{What is checked at the top (or where is the \texttt{break})?} and (2) \enquote{What changes inside the loop, so that check can change?}
In the homework (Slime Battle) the same bug happens if \texttt{health} never goes down.
\end{bughunt}

\begin{freebuild}[My own pixel picture]
Last 6 minutes: copy \texttt{INVADER}, give it a new name (e.g. \texttt{HEART}, \texttt{MUSHROOM}, \texttt{ROBOT}) and change the letters.
Design it on squared paper first if that helps. Every row must be the same length. Remember to change the line
\texttt{draw\_grid(t, INVADER, 30)} to use the new name.
\end{freebuild}

\section{Wrap-up (52--60 min)}
\begin{safety}[Fair play]
Ask: \enquote{Could we make the guessing game cheat? How?} (Change the secret number after each guess; say \enquote{too low} when it is not.)
\enquote{Would that be fun?} Then talk about tricks real games sometimes use: a timer that pressures you to buy something,
\enquote{only 1 left!} messages that are not true, rewards that make you feel you must come back every day, or nagging you to spend money.
Agree a class rule for our games: \textbf{the rules are honest, the player can stop any time, and nobody is pressured}.
Students can check their own games with this rule from now on.
\end{safety}

\begin{formativecheck}[Exit ticket]
How many times does this print \texttt{"jump"}, and what is the bug if we delete the last line?
\begin{pycode}
energy = 6
while energy > 0:
    print("jump")
    energy = energy - 2
\end{pycode}
\textbf{Answer:} 3 times (energy 6, 4, 2; then 0 stops it). Without the last line, energy stays 6 and it prints \texttt{jump} for ever: an infinite loop.
\textbf{Look for:} students who answer 2 or 4 need a trace table next week (the week 6 warm-up practices exactly this).
\end{formativecheck}

\begin{reflection}[Say it or write it]
\begin{itemize}[leftmargin=1.2em]
  \item When would you choose \texttt{while} instead of \texttt{for}?
  \item What did you do when your program got stuck in a loop?
  \item In the pixel art, which loop is the outer one, and what does it do?
\end{itemize}
\end{reflection}
Hand out homework 4 (Slime Battle). It takes 15--20 minutes.

\section{Notes for the teacher}
\begin{commonmistakes}
  \item \textbf{Nothing changes inside the loop.} The condition can never become false. Fix: change the variable inside the loop (\texttt{n = n - 1}).
  \item \textbf{The change is outside the loop.} \texttt{tries = tries + 1} written with no indent runs once, after the loop. Indentation decides what is inside.
  \item \textbf{\texttt{int()} before checking.} \texttt{int("banana")} crashes. Check with \texttt{answer.isdigit()} first, then convert.
  \item \textbf{\texttt{break} in the wrong place.} A \texttt{break} that is not inside an \texttt{if} leaves the loop on the first round.
  \item \textbf{Pixel art: inner loop not indented under the outer loop}, or \texttt{x} not reset at the start of each row, so the picture slants off to the right.
  \item \textbf{Rows of different lengths} make a lopsided picture. Count the characters in each row.
\end{commonmistakes}

\begin{differentiation}
\supportpath Do the countdown trace on paper together before touching the computer. In 4.1, give the pair the skeleton for TODO 2 on paper:
\texttt{while True:} / \texttt{answer = input(...)} / \texttt{if answer.isdigit():} / \texttt{return int(answer)} and add the range check later.
Pixel Art can move to next week's free-build time.
\extensionpath Limit the guessing game to 5 tries (\texttt{while tries < 5}) and print \enquote{Out of tries!} with the secret number.
In Pixel Art, make \texttt{draw\_grid} take a starting \texttt{x} and \texttt{y} and draw the invader three times in a row.
\end{differentiation}

\begin{stretchchallenge}[Computer guesses YOUR number]
Turn the game round: the player thinks of a number from 1 to 100 and the computer guesses, using the halving trick from the warm-up.
Keep \texttt{low = 1} and \texttt{high = 100}; guess \texttt{(low + high) // 2}; if the player says \enquote{higher}, set \texttt{low = guess + 1},
if \enquote{lower}, \texttt{high = guess - 1}. How many guesses does it ever need? (At most 7.)
\end{stretchchallenge}

\begin{teachernote}[If you are short on time]
Shorten the warm-up to one round. Do Exercise 4.1 fully (it holds the infinite-loop bug hunt and input validation) and give Pixel Art
as an extra at home or in buffer session B2. Keep the fair-play talk: it takes 3 minutes.
\end{teachernote}
